% **************************************************************************
%           Codeanalyse + Ergebnisanzeige
% **************************************************************************
\section{Tools zur Visualisierung von Testdaten zur Sicherung der Software"=Qualität}
\label{ch:SdF:Tools}

In diesem Abschnitt werden die beiden Tools Sonar und Grafana vorgestellt, welche im späteren Projektkontext Anwendung finden. Beide Tools ermöglichen eine Dash\-board"=Ansicht zur Darstellung erreichter Testergebnisse, um ein Monitoring der Software"=Qualität zu ermöglichen.

% ***********************************************************************
\subsection{Codeanalysetool Sonar}
\label{ch:SdF:Tools:sec:Sonar}

Sonar ist ein Werkzeug, um den Zustand und die Qualität von Software"=Code zu analysieren und deckt dabei viele Programmiersprachen ab. Anhand der eigenen Testsuite stellt Sonar dessen Ergebnisse, wie beispielsweise die erreichte Code"=Coverage (dt. Code"=Abdeckung), in einer Dashboard"=Übersicht dar. Enstprechend der individuellen Anwendungsfälle, kann Sonar in verschiedenen Editionen erworben werden. Dabei wird neben einer selbstgehosteten Instanz mit SonarQube eine Cloud"=Lösung angeboten. Mit SonarLint wird weiterhin die Integration eines Linters zum Verbessern der Code"=Qualität angeboten
\footnote{Sonar. Sonarqube. URL: \url{https://www.sonarsource.com/sem/products/sonarqube\} [Zugegriffen am 03.09.2025]}
\footnote{Sonar. Plans and Pricing. URL: \url{https://www.sonarsource.com/de/plans-and-pricing/} [Zugegriffen am 03.09.2025]}
\footnote{Sonar. Linter IDE Tool und Echtzeitsoftware für Code mit SonarQube for IDE. URL: \url{https://www.sonarsource.com/de/products/sonarlint/} [Zugegriffen am 03.09.2025]}.

% ***********************************************************************
\subsection{Dashboard-Framework Grafana}
\label{ch:SdF:Tool:sec:Grafana}

Das Dashboard-Framework Grafana ermöglicht es, Daten graphisch und übersichtlich in Form eines Dash\-boards (dt. Anzeigetafel) darzustellen. Dabei ermöglichen verschiedene Open Source Plugins (dt. Plugins aus offener Quelle), viele Daten"=Formate auszulesen und diese an die Dash\-board"=Elemente über einen API-Zugang zu übertragen und visualisieren. Mit einer großen Auswahl an Visualisierungsmöglichkeiten, ermöglicht Grafana, Daten angepasst auf individuelle Bedürnisse darzustellen. Je nach Nutzung bietet Grafana eine Enterprise"= und Open Source Version an, wodurch es in der Cloud als auch selbstgehostet eingesetzt werden kann.
\footnote{Grafana Labs. Dashboards. URL: \url{https://grafana.com/docs/grafana/latest/dashboards/} [Zugegriffen am 03.09.2025]}
\footnote{Michelle Tan. Grafana Labs. Celebrating Grafana 10: Top 1 Grafana features you need to know about. 2023. URL: \url{https://grafana.com/blog/2023/07/14/celebrating-grafana-10-top-10-grafana-features-you-need-to-know-about/} [Zugegriffen am 03.09.2025]}
\footnote{Grafana Labs. Build Dashboards. URL: \url{https://grafana.com/docs/grafana/latest/dashboards/build-dashboards/} [Zugegriffen am 03.09.2025]}
